% !TEX program = xelatex
% IMPORTANTE: compile com XeLaTeX (ou LuaLaTeX), NÃO pdfLaTeX — o template usa
% as fontes Arial/Times via fontspec. No Overleaf: Menu > Settings > Compiler > XeLaTeX.
% =====================================================================
% TEMPLATE DE TRABALHO ACADÊMICO (Dissertação de Mestrado / Tese de Doutorado)
% Programa de Pós-Graduação — UNIVALI
% Em conformidade com:
%   ABNT NBR 14724:2024 (apresentação de trabalhos acadêmicos)
%   ABNT NBR 6023:2025   (referências)
%   ABNT NBR 10520:2023  (citações)
% ---------------------------------------------------------------------
% COMO USAR (resumo — detalhes no README.md):
%   1. Preencha os METADADOS abaixo (\titulo..., \autor, \orientador...).
%   2. Escreva os capítulos em tex/*.tex (substitua o texto de exemplo).
%   3. Cadastre suas referências em referencias.bib.
%   4. Compile com XeLaTeX ou LuaLaTeX + biber (ver README.md).
% Os textos abaixo estão em "lorem ipsum"/placeholder: substitua-os.
% =====================================================================

% --- OPÇÕES DA CLASSE -------------------------------------------------
% NÍVEL DO TRABALHO (escolha UM):
%   graduacao | tcc        = "Trabalho de Conclusão de Curso", grau Bacharel
%   especializacao | monografia = "Monografia", grau Especialista
%   mestrado               = "Dissertação", grau Mestre/Mestra
%   doutorado              = "Tese", grau Doutor/Doutora  (padrão)
% MODIFICADORES (opcionais):
%   qualificacao = versão de qualificação (só mestrado/doutorado; suprime carimbo)
%   pre-defesa   = carimba "EXEMPLAR DE DEFESA"
%   pos-defesa   = carimba "VERSÃO FINAL" (padrão)
%   impressao    = versão para gráfica (desativa capa/contracapa no PDF)
% Exemplos:
%   [mestrado,qualificacao]   [doutorado,pos-defesa]
%   [tcc]                     [especializacao]
\documentclass[mestrado,qualificacao]{packages/ppg}

\usepackage[brazil]{babel}
\usepackage{csquotes}            % recomendado para biblatex + babel

% Arquivo de referências (biblatex-abnt)
\addbibresource{referencias.bib}

% --- PACOTES OPCIONAIS (mantenha apenas os que usar) ------------------
\usepackage{multirow}            % células mescladas em tabelas
\usepackage{makecell}            % \thead / \headrow em cabeçalhos
\usepackage{booktabs}            % \toprule \midrule \bottomrule
\usepackage{longtable}           % tabelas que cruzam páginas
\usepackage{tikz}                % desenhos e diagramas
\usetikzlibrary{shapes.geometric,arrows.meta,positioning}
\usepackage{pgfplots}            % gráficos (barras, dispersão, linhas)
\pgfplotsset{compat=1.17}
\usepackage{enumitem}            % listas com rótulos personalizados

\addto\captionsbrazil{\renewcommand{\listfigurename}{Lista de figuras}}

% --- ESTILO PARA LISTAGENS DE CÓDIGO (listings; opcional) -------------
\definecolor{codegreen}{RGB}{64, 180, 64}
\definecolor{codegray}{RGB}{130, 130, 130}
\definecolor{backcolour}{RGB}{245, 245, 235}
\lstdefinestyle{code}{
    backgroundcolor = \color{backcolour},
    commentstyle    = \color{codegreen},
    keywordstyle    = \color{magenta},
    numberstyle     = \tiny\color{codegray},
    basicstyle      = \ttfamily\scriptsize,
    breaklines      = true,
    captionpos      = t,   % legenda ACIMA do código (NBR 14724:2024, 5.8)
    numbers         = left,
    numbersep       = 5pt,
    tabsize         = 2,
}
\lstset{style=code}

% Utilitários de placeholder (caixa de figura). Remova quando não precisar mais.
% =====================================================================
% Utilitários de PLACEHOLDER para o template.
% Pode remover este arquivo (e o \input dele em thesis.tex) quando o
% trabalho estiver com figuras reais.
% =====================================================================

% Caixa cinza tracejada que ocupa o lugar de uma figura ainda não inserida.
% Uso: \figuraplaceholder{largura}{altura}{caminho-sugerido}
%   ex.: \figuraplaceholder{0.6\textwidth}{4cm}{imagens/sua-figura.png}
% Para inserir a imagem real, troque a chamada por:
%   \includegraphics[width=0.6\textwidth]{imagens/sua-figura.png}
\newcommand{\figuraplaceholder}[3]{%
  \begin{tikzpicture}
    \node[draw=gray, fill=gray!12, thick, dashed, rounded corners,
          minimum width=#1, minimum height=#2, align=center,
          text=gray!60, font=\small, inner sep=6pt]
      {[ FIGURA --- substitua por\\ \texttt{\textbackslash includegraphics\{#3\}} ]};
  \end{tikzpicture}%
}


% =====================================================================
% METADADOS DA CAPA E FOLHA DE ROSTO  — PREENCHA AQUI
% Na capa/folha de rosto, apenas a 1ª letra (ou nomes próprios) em maiúscula.
% =====================================================================
\tituloPT{Título do trabalho em português: subtítulo, se houver}
\tituloEN{Title of the work in English: subtitle, if any}
\autor[Sobrenome, N. N.]{Nome Completo do Autor}
\genero{M}                       % M = Masculino / F = Feminino
\orientador[Orientador]{Prof. Dr.}{Nome do Orientador}
\coorientador[Coorientador]{Prof. Dr.}{Nome do Coorientador} % opcional
\data{31}{12}{2026}              % data do depósito (dia, mês, ano)
\idioma{PT}                      % idioma principal (PT ou EN)
\local{Itajaí}                   % cidade na capa e folha de rosto

% --- INSTITUIÇÃO / CURSO ----------------------------------------------
% OPÇÃO A — Pós-graduação UNIVALI pré-configurada (PPGCA): basta a sigla.
% Já define área e especialidade automaticamente.
\curso{PPGC}
%
% OPÇÃO B — QUALQUER OUTRO curso (inclusive TCC/monografia): comente a linha
% acima, descomente as três linhas abaixo e preencha. Aqui você define a área
% e a especialidade manualmente (não use \curso{PPGC} junto, pois ele as
% sobrescreve).
%\curso{Curso de Ciência da Computação -- UNIVALI}   % nome do curso/instituição
%\area{Engenharia de Software}                        % aparece na folha de rosto
%\especialidade{em Ciência da Computação}             % completa "...título de <grau> <especialidade>"

% Cabeçalho da instituição no TOPO DA CAPA (opcional; só redefina se NÃO usar o PPGCA).
% Útil em TCC/monografia. Use \\ para quebrar linhas; será exibido em maiúsculas.
%\capainstituicao{UNIVERSIDADE DO VALE DO ITAJAÍ\\[6pt] ~\\[6pt] CENTRO DE CIÊNCIAS TECNOLÓGICAS\\ CURSO DE CIÊNCIA DA COMPUTAÇÃO}

% =====================================================================
% RESUMO (PORTUGUÊS) — máx. 500 palavras, de 1 a 5 palavras-chave
% =====================================================================
\textoresumo[brazil]{
Lorem ipsum dolor sit amet, consectetur adipiscing elit. Sed do eiusmod tempor
incididunt ut labore et dolore magna aliqua. Ut enim ad minim veniam, quis
nostrud exercitation ullamco laboris nisi ut aliquip ex ea commodo consequat.
Duis aute irure dolor in reprehenderit in voluptate velit esse cillum dolore eu
fugiat nulla pariatur. Excepteur sint occaecat cupidatat non proident, sunt in
culpa qui officia deserunt mollit anim id est laborum. [Substitua por: contexto,
problema, objetivo, método, principais resultados e conclusão do seu trabalho.]
}
{Palavra-chave um. Palavra-chave dois. Palavra-chave três.}

% =====================================================================
% RESUMO (INGLÊS / ABSTRACT) — máx. 500 palavras, de 1 a 5 palavras-chave
% =====================================================================
\textoresumo[english]{
Lorem ipsum dolor sit amet, consectetur adipiscing elit, sed do eiusmod tempor
incididunt ut labore et dolore magna aliqua. Ut enim ad minim veniam, quis
nostrud exercitation ullamco laboris nisi ut aliquip ex ea commodo consequat.
[Replace with the English version of your abstract: context, problem, objective,
method, main results, and conclusion.]
}
{Keyword one. Keyword two. Keyword three.}

% =====================================================================
% ELEMENTOS PRÉ-TEXTUAIS
% =====================================================================
% Ficha catalográfica — OPCIONAL na qualificação; obrigatória na versão final
% (elaborada por bibliotecário(a)). Para incluir, descomente:
%\incluifichacatalografica{tex/pre-textual/ficha-catalografica.tex}

% Folha de aprovação — obrigatória na versão FINAL (após a defesa). Descomente:
%\incluifolhadeaprovacao{tex/pre-textual/folha-aprovacao.tex}

% Dedicatória / Agradecimentos / Epígrafe — todos OPCIONAIS.
\textodedicatoria*{tex/pre-textual/dedicatoria}
\textoagradecimentos*{tex/pre-textual/agradecimentos}
\textoepigrafe*{tex/pre-textual/epigrafe}

% Listas (descomente as que for usar)
\incluilistadefiguras
%\incluilistadetabelas
\incluilistadequadros
%\incluilistadealgoritmos
%\incluilistadecodigos
\incluilistadesiglas

% =====================================================================
% LISTA DE ABREVIATURAS (grupo A) E SIGLAS (grupo S)
% Cadastre aqui cada sigla com \nomenclature; ela aparece automaticamente
% na "Lista de siglas". Exemplos abaixo — substitua pelos seus.
% =====================================================================
\nomenclature[A]{ed.}{Edição}
\nomenclature[A]{p.}{Página}

\nomenclature[S]{ABNT}{Associação Brasileira de Normas Técnicas}
\nomenclature[S]{UNIVALI}{Universidade do Vale do Itajaí}
\nomenclature[S]{PPG}{Programa de Pós-Graduação}

% =====================================================================
% INÍCIO DO DOCUMENTO
% =====================================================================
\begin{document}

% Declaração de uso de IA Generativa (Política CNPq, Portaria 2.664/2026, Art. 9).
% Remova esta linha se não tiver utilizado IA generativa no trabalho.
\chapter*{Declaração de Uso de Inteligência Artificial Generativa}

% Em conformidade com o Art. 9º da Política de Integridade na Atividade
% Científica do CNPq (Portaria nº 2.664/2026). Ajuste conforme o seu uso,
% ou remova o \input deste arquivo em thesis.tex se não tiver usado IA.

Declara-se o uso da(s) ferramenta(s) de inteligência artificial generativa
[nome(s) da(s) ferramenta(s)] no desenvolvimento deste trabalho, com as
seguintes finalidades: [ex.: apoio à redação e estruturação do texto; revisão
linguística; organização de referências; geração de código]. O autor revisou,
verificou e validou todo o conteúdo, sendo integralmente responsável pelo texto
final, pelos dados e pelas referências apresentados.


% ----------------------------------------------------------
% ELEMENTOS TEXTUAIS
% ----------------------------------------------------------
\textual

\chapter{Introdução}
\label{chp:introducao}
% =====================================================================
% INTRODUÇÃO
% Apresente: tema e contexto, problema, objetivos, justificativa e a
% organização do texto. Substitua todo o lorem ipsum pelo seu conteúdo.
% Dica de citação (biblatex-abnt):
%   \parencite{chave}      -> (SOBRENOME, ano)
%   \textcite{chave}       -> Sobrenome (ano)
%   \parencite[p. 12]{chave} -> (SOBRENOME, ano, p. 12)
% =====================================================================

Lorem ipsum dolor sit amet, consectetur adipiscing elit, sed do eiusmod tempor
incididunt ut labore et dolore magna aliqua~\parencite{exemplo-livro}. Ut enim ad
minim veniam, quis nostrud exercitation ullamco laboris nisi ut aliquip ex ea
commodo consequat. Segundo \textcite{exemplo-artigo}, duis aute irure dolor in
reprehenderit in voluptate velit esse cillum dolore eu fugiat nulla pariatur.

Excepteur sint occaecat cupidatat non proident, sunt in culpa qui officia
deserunt mollit anim id est laborum~\parencite{exemplo-artigo, exemplo-livro}. Sed
ut perspiciatis unde omnis iste natus error sit voluptatem accusantium doloremque
laudantium, totam rem aperiam, eaque ipsa quae ab illo inventore veritatis.

%-------------------------------------------------------------
\section{Problema de pesquisa}
%-------------------------------------------------------------

Nemo enim ipsam voluptatem quia voluptas sit aspernatur aut odit aut fugit, sed
quia consequuntur magni dolores eos qui ratione voluptatem sequi nesciunt. A
questão que orienta este trabalho pode ser sintetizada como:

\begin{itemize}
  \item Qual é a primeira questão de pesquisa? Lorem ipsum dolor sit amet.
  \item Qual é a segunda questão de pesquisa? Consectetur adipiscing elit.
\end{itemize}

%-------------------------------------------------------------
\section{Objetivos}
%-------------------------------------------------------------

Nesta seção apresentam-se o objetivo geral e os objetivos específicos.

\subsection{Objetivo geral}

Neque porro quisquam est, qui dolorem ipsum quia dolor sit amet, consectetur,
adipisci velit, sed quia non numquam eius modi tempora incidunt ut labore.

\subsection{Objetivos específicos}

\begin{itemize}
  \item Primeiro objetivo específico: ut enim ad minima veniam, quis nostrum.
  \item Segundo objetivo específico: exercitationem ullam corporis suscipit.
  \item Terceiro objetivo específico: laboriosam, nisi ut aliquid ex ea commodi.
\end{itemize}

%-------------------------------------------------------------
\section{Justificativa}
%-------------------------------------------------------------

Quis autem vel eum iure reprehenderit qui in ea voluptate velit esse quam nihil
molestiae consequatur, vel illum qui dolorem eum fugiat quo voluptas nulla
pariatur~\parencite{exemplo-livro}. At vero eos et accusamus et iusto odio
dignissimos ducimus qui blanditiis praesentium voluptatum deleniti atque.

%-------------------------------------------------------------
\section{Estrutura do trabalho}
%-------------------------------------------------------------

Este trabalho está organizado em cinco capítulos. O Capítulo~\ref{chp:introducao}
apresentou a introdução. O Capítulo~\ref{chp:fundamentacao-teorica} reúne a
fundamentação teórica. O Capítulo~\ref{chp:materiais-metodos} descreve os
materiais e métodos. O Capítulo~\ref{chp:resultados} apresenta os resultados e a
discussão. Por fim, o Capítulo~\ref{chp:consideracoes-finais} traz as
considerações finais.


\chapter{Fundamentação Teórica}
\label{chp:fundamentacao-teorica}
% =====================================================================
% FUNDAMENTAÇÃO TEÓRICA (Referencial Teórico)
% Reúna os conceitos necessários para o leitor compreender o trabalho.
% Este capítulo demonstra: figura, equação, citação e listas.
% =====================================================================

Lorem ipsum dolor sit amet, consectetur adipiscing elit, sed do eiusmod tempor
incididunt ut labore et dolore magna aliqua~\parencite{exemplo-livro}. Ut enim ad
minim veniam, quis nostrud exercitation ullamco laboris.

%-------------------------------------------------------------
\section{Primeiro conceito}
\label{sec:conceito-um}
%-------------------------------------------------------------

Duis aute irure dolor in reprehenderit in voluptate velit esse cillum dolore eu
fugiat nulla pariatur~\parencite{exemplo-artigo}. A \autoref{fig:exemplo} ilustra
o conceito.

% Citação DIRETA curta (até 3 linhas): fica no corpo do texto, entre aspas, com a
% fonte e a página logo após (NBR 10520:2023). Use \enquote{...} (csquotes) para as
% aspas — NUNCA a aspa reta ", que sob XeLaTeX sai como ”...” (fecha nas duas pontas).
Nesse sentido, \enquote{a citação direta de até três linhas permanece no corpo do
texto, entre aspas}~\parencite[p. 15]{exemplo-livro}.

% Citação DIRETA LONGA (mais de 3 linhas): destacada em parágrafo próprio, com
% recuo de 4 cm da margem esquerda, fonte menor, espaço simples e SEM aspas
% (NBR 10520:2023; guia UNIVALI). Use o ambiente \begin{citacao}...\end{citacao}:
% o abnTeX2 já aplica recuo (4 cm), \ABNTEXfontereduzida e espaçamento simples.
% A chamada (autor, ano, página) precede o bloco; não se usam aspas.
Já a citação com mais de três linhas é destacada do texto, como em
\textcite[p. 86]{exemplo-livro}:

\begin{citacao}
  Lorem ipsum dolor sit amet, consectetur adipiscing elit, sed do eiusmod tempor
  incididunt ut labore et dolore magna aliqua. Ut enim ad minim veniam, quis
  nostrud exercitation ullamco laboris nisi ut aliquip ex ea commodo consequat.
  Duis aute irure dolor in reprehenderit in voluptate velit esse cillum dolore eu
  fugiat nulla pariatur.
\end{citacao}

% --- EXEMPLO DE FIGURA -----------------------------------------------
% Para usar uma imagem real, troque o conteúdo do \centering pela linha:
%   \includegraphics[width=0.75\textwidth]{imagens/sua-figura.png}
% O comando \figuraplaceholder (definido em tex/_placeholders.tex) apenas
% desenha uma caixa cinza para o template compilar sem imagens externas.
\begin{figure}[!htb]
  \caption{Legenda da figura (acima da figura, conforme NBR 14724).}
  \label{fig:exemplo}
  \centering
  % A fonte fica alinhada à margem esquerda DA FIGURA (não do texto).
  \elementocomfonte{\figuraplaceholder{0.6\textwidth}{4cm}{imagens/sua-figura.png}}%
    {Fonte: Elaborada pelo autor (\the\year{}).}
\end{figure}

%-------------------------------------------------------------
\section{Formulação matemática}
\label{sec:formulacao}
%-------------------------------------------------------------

Equações são numeradas e podem ser referenciadas, como mostra a
\autoref{eq:exemplo}:

\begin{equation}
  E = m c^{2},
  \label{eq:exemplo}
\end{equation}

\noindent onde $E$ é a energia, $m$ a massa e $c$ a velocidade da luz no vácuo.
At vero eos et accusamus et iusto odio dignissimos ducimus.

%-------------------------------------------------------------
\section{Síntese}
%-------------------------------------------------------------

Os principais pontos abordados foram:

\begin{enumerate}[label=(\roman*)]
  \item primeiro ponto: lorem ipsum dolor sit amet;
  \item segundo ponto: consectetur adipiscing elit;
  \item terceiro ponto: sed do eiusmod tempor incididunt.
\end{enumerate}


\chapter{Materiais e Métodos}
\label{chp:materiais-metodos}
% =====================================================================
% MATERIAIS E MÉTODOS (Metodologia)
% Descreva o tipo de pesquisa, materiais, procedimentos e instrumentos.
% Este capítulo demonstra: quadro (tabela com borda), tabela booktabs e
% listagem de código.
% =====================================================================

Lorem ipsum dolor sit amet, consectetur adipiscing elit. Este trabalho adota
abordagem qualitativa, quantitativa ou mista e caracteriza-se como pesquisa
exploratória, descritiva ou explicativa~\parencite{exemplo-livro}. [Descreva o
tipo de pesquisa adotado.]

%-------------------------------------------------------------
\section{Materiais}
\label{sec:materiais}
%-------------------------------------------------------------

O \autoref{quad:materiais} lista os materiais utilizados.

% --- EXEMPLO DE QUADRO (tabela com bordas; legenda acima) ------------
\begin{quadro}[!htb]
  \caption{Materiais e instrumentos utilizados.}
  \label{quad:materiais}
  \centering
  % A fonte fica alinhada à margem esquerda DO QUADRO (não do texto).
  \elementocomfonte{%
  \begin{tabular}{|l|l|c|}
    \hline
    \textbf{Item} & \textbf{Descrição} & \textbf{Quantidade} \\
    \hline
    Item A & Lorem ipsum dolor sit amet & 2 \\
    \hline
    Item B & Consectetur adipiscing elit & 5 \\
    \hline
    Item C & Sed do eiusmod tempor        & 1 \\
    \hline
  \end{tabular}}%
  {Fonte: Elaborado pelo autor (\the\year{}).}
\end{quadro}

A \autoref{tab:dados} apresenta os mesmos dados no estilo \textit{booktabs}
(tabela aberta, sem bordas verticais), recomendado para dados numéricos.

% --- EXEMPLO DE TABELA (estilo booktabs) -----------------------------
\begin{table}[!htb]
  \caption{Exemplo de tabela numérica.}
  \label{tab:dados}
  \centering
  % A fonte fica alinhada à margem esquerda DA TABELA (não do texto).
  \elementocomfonte{%
  \begin{tabular}{lrr}
    \toprule
    \textbf{Categoria} & \textbf{Valor} & \textbf{\%} \\
    \midrule
    Primeira  & 120 & 48,0 \\
    Segunda   &  90 & 36,0 \\
    Terceira  &  40 & 16,0 \\
    \midrule
    \textbf{Total} & \textbf{250} & \textbf{100,0} \\
    \bottomrule
  \end{tabular}}%
  {Fonte: Elaborada pelo autor (\the\year{}).}
\end{table}

%-------------------------------------------------------------
\section{Procedimentos}
\label{sec:procedimentos}
%-------------------------------------------------------------

Os procedimentos foram organizados nas seguintes etapas:

\begin{enumerate}[label=\arabic*.]
  \item \textbf{Etapa um:} lorem ipsum dolor sit amet, consectetur adipiscing;
  \item \textbf{Etapa dois:} sed do eiusmod tempor incididunt ut labore;
  \item \textbf{Etapa três:} ut enim ad minim veniam, quis nostrud exercitation.
\end{enumerate}

O \autoref{lst:exemplo} mostra um exemplo de trecho de código.

% --- EXEMPLO DE LISTAGEM DE CÓDIGO -----------------------------------
\begin{lstlisting}[language=Python, caption={Exemplo de código.}, label={lst:exemplo}]
def media(valores):
    # Calcula a media aritmetica de uma lista
    return sum(valores) / len(valores)

print(media([120, 90, 40]))
\end{lstlisting}
% Fonte do código, alinhada à esquerda e em tamanho reduzido (uniforme com a
% legenda), conforme NBR 14724:2024 (5.8).
\noindent{\ABNTEXfontereduzida Fonte: Elaborado pelo autor (\the\year{}).}


\chapter{Resultados e Discussão}
\label{chp:resultados}
% =====================================================================
% RESULTADOS E DISCUSSÃO
% Apresente os resultados (texto, tabelas, gráficos) e discuta-os à luz
% da fundamentação teórica. Este capítulo demonstra um gráfico com pgfplots.
% =====================================================================

Lorem ipsum dolor sit amet, consectetur adipiscing elit, sed do eiusmod tempor
incididunt ut labore et dolore magna aliqua. Os resultados são apresentados na
\autoref{fig:grafico} e discutidos a seguir.

%-------------------------------------------------------------
\section{Resultados}
\label{sec:resultados}
%-------------------------------------------------------------

% --- EXEMPLO DE GRÁFICO (pgfplots, autocontido) ----------------------
\begin{figure}[!htb]
  \caption{Exemplo de gráfico de barras gerado com pgfplots.}
  \label{fig:grafico}
  \centering
  % A fonte fica alinhada à margem esquerda DO GRÁFICO (não do texto).
  \elementocomfonte{%
  \begin{tikzpicture}
  \begin{axis}[
    ybar, width=0.85\textwidth, height=6cm,
    symbolic x coords={Primeira,Segunda,Terceira},
    xtick=data,
    ylabel={Valor (\%)},
    ymin=0, ymax=60,
    bar width=24pt,
    nodes near coords, nodes near coords style={font=\footnotesize},
    enlarge x limits=0.25,
  ]
  \addplot coordinates {(Primeira,48) (Segunda,36) (Terceira,16)};
  \end{axis}
  \end{tikzpicture}}%
  {Fonte: Elaborada pelo autor (\the\year{}).}
\end{figure}

Ut enim ad minim veniam, quis nostrud exercitation ullamco laboris nisi ut
aliquip ex ea commodo consequat~\parencite{exemplo-artigo}.

%-------------------------------------------------------------
\section{Discussão}
\label{sec:discussao}
%-------------------------------------------------------------

Duis aute irure dolor in reprehenderit in voluptate velit esse cillum dolore eu
fugiat nulla pariatur. Os achados são coerentes com \textcite{exemplo-livro} e
contrastam com parte da literatura, sugerindo que [interprete aqui os seus
resultados].


\chapter{Considerações Finais}
\label{chp:consideracoes-finais}
% =====================================================================
% CONSIDERAÇÕES FINAIS (Conclusão)
% Retome o problema e os objetivos, sintetize os resultados, aponte
% limitações e trabalhos futuros. NÃO introduza conteúdo novo aqui.
% =====================================================================

Lorem ipsum dolor sit amet, consectetur adipiscing elit, sed do eiusmod tempor
incididunt ut labore et dolore magna aliqua. Este trabalho teve por objetivo
[retome o objetivo geral] e os resultados indicam que [síntese dos achados].

Quanto aos objetivos específicos, conclui-se que: (i) lorem ipsum dolor sit amet;
(ii) consectetur adipiscing elit; e (iii) sed do eiusmod tempor incididunt.

%-------------------------------------------------------------
\section{Limitações}
%-------------------------------------------------------------

Ut enim ad minim veniam, quis nostrud exercitation ullamco laboris nisi ut
aliquip ex ea commodo consequat.

%-------------------------------------------------------------
\section{Trabalhos futuros}
%-------------------------------------------------------------

\begin{itemize}
  \item Primeira direção de trabalho futuro: lorem ipsum dolor sit amet;
  \item Segunda direção de trabalho futuro: consectetur adipiscing elit.
\end{itemize}


% ----------------------------------------------------------
% ELEMENTOS PÓS-TEXTUAIS
% ----------------------------------------------------------
\bookmarksetup{startatroot}
\postextual

% Lista de siglas/abreviaturas (glossário) e Referências
\printglossary
\printbibliography[title={REFERÊNCIAS}]

% ----------------------------------------------------------
% APÊNDICES (produção do próprio autor) — OPCIONAL
% ----------------------------------------------------------
%\begin{apendicesenv}
%    \chapter{Título do apêndice}
%    \label{ape:exemplo}
%    Conteúdo do apêndice.
%\end{apendicesenv}

% ----------------------------------------------------------
% ANEXOS (material de terceiros) — OPCIONAL
% ----------------------------------------------------------
%\begin{anexosenv}
%    \chapter{Título do anexo}
%    \label{anx:exemplo}
%    Conteúdo do anexo.
%\end{anexosenv}

\end{document}
